\documentclass[12pt]{article}
\usepackage[utf8]{inputenc}
\usepackage[T1]{fontenc}
\usepackage{hyperref}
\title{Verre Language Resources for K-12 Curriculum Development}
\author{Compiled by Claude Sonnet 5 (claude-sonnet-5) for Andr\'as Kornai}
\date{}
\begin{document}
\maketitle

\section{Introduction}
Verre (Vere; also Kobo, Were, Ziri) is spoken in Adamawa State, Nigeria, and in Faro department, Cameroon. The name covers a cluster of related lects, such as Gweri and Bai, and Gweri was chosen as the standard for the Verre Language Project. Speaker figures differ by source. Little exists online.

\section{Primary Educational Resources}
A 2013 pamphlet, \emph{Reading and Writing Verre}, was produced by the Gweri Language Committee and the Luke Partnership. It is reported in search results, but no URL for it was seen.

\section{Linguistic and Reference Materials}
\begin{itemize}
\item Verre glossary in \emph{Vestiges} journal: \url{https://vestiges-journal.info/2021/Verre/HTML/15_app-3_verre-glossary.htm}
\item Kleinewillinghöfer, Vere, in Glottolog references: \url{https://glottolog.org/resource/reference/id/550010}
\item Overview: \url{https://en.wikipedia.org/wiki/Vere_language}
\end{itemize}

\section{Religious and Community Resources}
Joshua Project describes the Verre people: \url{https://joshuaproject.net/index.php/languages/ver}

\section{Digital and Audio Resources}
None found.

\section{Curriculum Development Context}
Verre is not known to be taught in schools. Only a literacy pamphlet and linguistic descriptions were found, so a curriculum would need new materials developed with the Gweri Language Committee.
\end{document}